\section{Discussion}
\label{sec:discussion}

\subsection{What these results imply for reporting}

The three parameters examined here are not exotic. Each is fixed in the course
of designing any evaluation of carbon-aware scheduling, and each is ordinarily
left unstated. Their combined effect, in our experiments, spans reported
reductions from 4.6\,\% to 90.7\,\% for the same class of policy. A figure
reported without them is therefore not comparable with a figure reported
elsewhere, and the divergence in the literature is to that extent expected
rather than anomalous.

Four items would make such figures comparable.

\emph{Capacity utilisation.} Total demand as a fraction of total available
slot-hours, stated alongside any claim about deadline flexibility. Without it, a
reduction attributed to flexibility cannot be distinguished from a reduction
attributable to an unconstrained pool.

\emph{Origin region.} For any policy that shifts in time but not in space, the
region in which the workload originates, and preferably the distribution of
results across candidate origins rather than a single value.

\emph{Region-set composition.} The candidate regions and their mean carbon
intensities, from which a reader can judge whether the set admits a dominant
optimum. The ratio between the cleanest region and the second cleanest is a
compact summary.

\emph{Accounting basis.} Whether reported emissions include the idle draw of
provisioned hosts, of active hosts only, or neither. As
Section~\ref{sec:accounting} shows, the first of these cannot distinguish
policies at all in a fixed-fleet setting, and the second conflates regional
placement with consolidation.

\subsection{Concentration and deployability}

The mechanism common to all three findings is concentration. Across every valid
cell in both region sets, the correlation between the share of work placed in
the busiest region and attributed emissions is below $-0.93$. The policies
reduce emissions by moving work into the cleanest reachable grid, and the three
parameters govern how much of the workload can be moved, where a constrained
policy starts from, and how much is gained by arriving.

This has an uncomfortable consequence. The configurations that report the
largest reductions are those in which 95\,\% or more of the workload is placed
in one region, and that is not a configuration any operator could adopt. Real
deployments are constrained by data residency, by latency to users, by capacity
that is finite in exactly the way our capped configurations model, and by the
grid-level effects of concentrating demand. The reductions reported at low
utilisation should therefore be read as upper bounds obtained under a placement
distribution that would not survive contact with those constraints, rather than
as achievable savings.

Read the other way, the result is more encouraging than it first appears. The
82\,\% utilisation row, where deadline flexibility is worth 0.9\,\% and space
shifting 4.6\,\%, is the row that most resembles an efficiently operated
facility. That the returns are small there is a finding about where effort is
best directed, not an argument against carbon-aware scheduling as such.

\subsection{On the interpretation of time shifting}

Section~\ref{sec:res-home} shows that a time-shifting policy's outcome tracks
the mean intensity of its origin rather than that origin's temporal structure.
The natural reading is that temporal flexibility is worth little. The more
precise reading is that temporal flexibility is worth little \emph{when spatial
flexibility is available and capacity is not binding}, which is the regime in
which the two are usually compared. Our delay-insensitive robustness condition
supplies the counterexample: at tight capacity and wide margin, where
relocation is constrained and deferral is not, the correlation with origin
intensity falls to $+0.403$ and that with diurnal amplitude rises to $+0.776$.
Temporal structure becomes the dominant term exactly when it is the only term
available.

A second observation from that condition bears on the premise of deadline-based
scheduling. The workload class the provider designates as delay-insensitive has
a mean lifetime an order of magnitude longer than the general mixture, so that a
given absolute margin represents a far smaller fraction of runtime. Deadline
margins expressed in hours are not comparable across workload classes, and
margins expressed as a fraction of runtime would be a more meaningful axis for
future studies.

\section{Limitations}
\label{sec:limitations}

\emph{Perfect foresight.} Planning is performed against the ground-truth carbon
trace. This isolates the parameters under study from prediction error, but it
means the absolute reductions reported here are upper bounds with respect to
forecast quality. Whether the three effects survive realistic forecast error is
a question we address separately.

\emph{Single workload source.} All virtual machine requests derive from one
provider's 2019 trace. The duration and size distributions of that trace shape
the capacity utilisations at which our effects appear, and a workload with a
different shape would place the transition points elsewhere, though we have no
reason to expect the direction of the effects to change.

\emph{Deadline margins are imposed rather than observed.} The trace records a
workload category but not a service-level objective, so deadline margins are
assigned as an experimental factor rather than read from the data. The
delay-insensitive condition is the closest available approximation to a
provider-designated flexibility class.

\emph{No migration or preemption.} Virtual machines run to completion where they
start, following the formulation of the work we build on. Policies permitting
mid-flight migration would have a larger action space, and the capacity
interaction we report might take a different form under them.

\emph{Excluded costs.} Network transfer between regions, the embodied carbon of
the hardware, and datacentre overhead beyond the host itself are not modelled.
The first of these is of particular relevance given that the largest reductions
we observe coincide with the greatest concentration of load.

\emph{Two region sets.} The claim of Section~\ref{sec:res-regionset} rests on a
contrast between two sets. The direction of the effect follows from the
mechanism and we expect it to generalise, but the coefficient relating
region-set composition to reported savings is not established by two points.
